%% Resumo
\begin{resumo}
A escassez de políticas públicas e ferramentas de visibilidade digital tem limitado o acesso e o reconhecimento de diversos pontos culturais, especialmente em regiões menos exploradas do Brasil. Esse contexto evidencia a necessidade de soluções que promovam o resgate e a valorização da cultura local, tornando-a mais acessível à população e aos visitantes.O uso de tecnologias digitais, em especial os aplicativos móveis, tem se mostrado uma estratégia viável e eficiente para solucionar essa demanda. Aplicativos com recursos como geolocalização, interface intuitiva e acessibilidade possibilitam registrar, divulgar e tornar interativos os espaços culturais e históricos, promovendo a democratização do acesso à cultura. Dessa forma, o presente trabalho propõe o desenvolvimento de um aplicativo de cadastro de pontos culturais, com funcionalidades voltadas à inclusão, navegação georreferenciada, apresentação de imagens e textos explicativos. A proposta busca contribuir para a valorização do patrimônio cultural do estado, ampliando sua visibilidade e facilitando o engajamento da população com sua própria história e identidade.

\vspace{\onelineskip}
\noindent
\textbf{Palavras-chaves}: Cultura Digital, Aplicativos Móveis, Geolocalização, Acessibilidade Cultural, Patrimônio Cultural, Inclusão Digital, Desenvolvimento Cultural, Tecnologias da Informação.
\end{resumo}